\documentclass[12pt,a4paper]{article}
\usepackage[utf8]{inputenc}
\usepackage[T1]{fontenc}
\usepackage[spanish]{babel}
\usepackage{geometry}
\usepackage{graphicx}
\usepackage{array}
\usepackage{tabularx}
\usepackage{booktabs}
\usepackage{xcolor}
\usepackage{enumitem}
\usepackage{hyperref}
\usepackage{fancyhdr}
\usepackage{titlesec}
\usepackage{listings}
\usepackage{float}
\usepackage{amsmath}

\geometry{margin=2.5cm}
\pagestyle{fancy}
\fancyhf{}
\fancyhead[L]{\small Electro+ Lab}
\fancyhead[R]{\small Documentación Técnica}
\fancyfoot[C]{\thepage}
\renewcommand{\headrulewidth}{0.4pt}

\definecolor{primary}{RGB}{31,77,58}
\definecolor{gold}{RGB}{201,168,106}
\definecolor{codebg}{RGB}{245,240,230}

\titleformat{\section}{\Large\bfseries\color{primary}}{}{0em}{}[\vspace{-0.3em}\rule{\textwidth}{0.5pt}]
\titleformat{\subsection}{\large\bfseries\color{primary!80!black}}{}{0em}{}

\lstset{
  basicstyle=\ttfamily\small,
  backgroundcolor=\color{codebg},
  frame=single,
  rulecolor=\color{gold},
  breaklines=true,
  columns=fullflexible,
  keepspaces=true,
}

\hypersetup{
  colorlinks=true,
  linkcolor=primary,
  urlcolor=primary,
  citecolor=primary,
}

\begin{document}

% ── Portada ──────────────────────────────────────────────────────────
\begin{titlepage}
  \centering
  \vspace*{2.5cm}
  {\Huge\bfseries\color{primary} Electro+ Lab\\[0.4em]}
  {\Large Simulador de Circuitos Eléctricos Interactivo\\[1.5cm]}
  {\large\textbf{Documentación Técnica para Evaluación Académica}\\[0.8cm]}
  {\normalsize Explicación del funcionamiento, arquitectura y construcción del sistema\\[2.5cm]}
  \begin{tabular}{rl}
    \textbf{Proyecto:} & LabCircuitos / Electro+ Lab \\
    \textbf{Tipo:} & Aplicación web educativa \\
    \textbf{Stack:} & React + TypeScript + Python + FastAPI \\
    \textbf{Simulación:} & MNA (Modified Nodal Analysis) \\
    \textbf{Fecha:} & \today \\
  \end{tabular}
  \vfill
  {\small Proyecto académico --- Uso educativo}
\end{titlepage}

\tableofcontents
\newpage

% ── 1. Introducción ──────────────────────────────────────────────────
\section{Introducción}

\textbf{Electro+ Lab} es un simulador de circuitos eléctricos que permite diseñar esquemas de forma visual, ejecutar simulaciones en tiempo real y visualizar mediciones mediante un multímetro virtual y un osciloscopio. El sistema está orientado al aprendizaje de conceptos fundamentales de ingeniería eléctrica: ley de Ohm, análisis nodal, comportamiento de semiconductores (LED, diodo), elementos reactivos (capacitor, inductor) e instrumentación de laboratorio.

\subsection{Objetivos del proyecto}

\begin{enumerate}[label=\arabic*.]
  \item Ofrecer un editor visual intuitivo para montar circuitos sobre una cuadrícula.
  \item Simular voltajes, corrientes y potencias con un motor numérico MNA en el navegador.
  \item Mostrar resultados en tiempo real en el canvas, multímetro y osciloscopio.
  \item Validar la topología del circuito antes de simular (tierra, cortocircuitos, nodos flotantes).
  \item Desplegar la aplicación en la nube (Vercel + Render) para acceso público.
\end{enumerate}

\subsection{Tecnologías utilizadas}

\begin{table}[H]
  \centering
  \begin{tabularx}{\textwidth}{lX}
    \toprule
    \textbf{Capa} & \textbf{Tecnologías} \\
    \midrule
    Frontend & React 18, TypeScript, Vite, Tailwind CSS, React Flow, Zustand, Plotly.js \\
    Motor de simulación & TypeScript puro (\texttt{src/engine/}), eliminación gaussiana, Euler hacia atrás \\
    Backend (opcional) & Python 3, FastAPI, NumPy, Uvicorn \\
    Despliegue & Vercel (frontend estático), Render (API) \\
    Herramientas & pnpm, ESLint, Prettier, Git/GitHub \\
    \bottomrule
  \end{tabularx}
  \caption{Stack tecnológico del proyecto}
\end{table}

% ── 2. Arquitectura ──────────────────────────────────────────────────
\section{Arquitectura del sistema}

El sistema sigue una arquitectura cliente-servidor \textbf{desacoplada}. La simulación principal ocurre en el \textbf{cliente} (navegador); el servidor es complementario.

\subsection{Diagrama de capas}

\begin{verbatim}
┌──────────────────────────────────────────────────────────────┐
│                    NAVEGADOR (Cliente)                       │
│  ┌─────────────┐  ┌──────────────┐  ┌───────────────────┐  │
│  │ Interfaz    │  │ Estado       │  │ Motor MNA         │  │
│  │ React       │──│ Zustand      │──│ src/engine/       │  │
│  │ (UI)        │  │ circuitStore │  │ TransientMNASolver│  │
│  └─────────────┘  └──────────────┘  └───────────────────┘  │
└──────────────────────────────────────────────────────────────┘
                         │ (opcional, HTTP)
                         ▼
┌──────────────────────────────────────────────────────────────┐
│              SERVIDOR Render (FastAPI + NumPy)               │
│  /api/simulate  ·  /api/simulate/validate  ·  /api/health  │
└──────────────────────────────────────────────────────────────┘
\end{verbatim}

\subsection{Separación de responsabilidades}

\begin{itemize}
  \item \textbf{Interfaz (\texttt{src/components/}, \texttt{src/features/}):} renderiza la UI, captura eventos del usuario y muestra mediciones. No contiene lógica de matrices.
  \item \textbf{Estado (\texttt{src/store/circuitStore.ts}):} almacena el circuito, cables, sondas, resultados y banderas de simulación.
  \item \textbf{Motor (\texttt{src/engine/}):} construye el grafo eléctrico, arma el sistema MNA y resuelve pasos transitorios. Sin dependencias de React.
  \item \textbf{Backend (\texttt{backend/}):} API REST para validación server-side, simulación de referencia y exportación SPICE.
\end{itemize}

% ── 3. Construcción del frontend ─────────────────────────────────────
\section{Cómo se construye el frontend}

\subsection{Punto de entrada}

\begin{itemize}
  \item \texttt{index.html} --- plantilla HTML que carga el bundle de Vite.
  \item \texttt{src/main.tsx} --- monta la aplicación React en el DOM.
  \item \texttt{src/App.tsx} --- layout principal con toolbar, canvas, paneles y osciloscopio.
\end{itemize}

\subsection{Editor visual (canvas)}

El editor usa la librería \textbf{React Flow} para representar nodos (componentes) y aristas (cables).

\begin{table}[H]
  \centering
  \small
  \begin{tabularx}{\textwidth}{lX}
    \toprule
    \textbf{Archivo} & \textbf{Función} \\
    \midrule
    \texttt{CircuitEditor.tsx} & Canvas con zoom, pan y cuadrícula de laboratorio \\
    \texttt{ComponentNode.tsx} & Nodo visual: símbolo SVG + lecturas V/I en vivo \\
    \texttt{symbols/index.tsx} & Símbolos eléctricos normalizados (IEC simplificado) \\
    \texttt{wireToEdge.ts} & Convierte cables del store en aristas de React Flow \\
    \texttt{componentHandles.ts} & Posición de bornes (+/−) según tipo de componente \\
    \texttt{constants.ts} & Cuadrícula (30 px), colores, plantillas de parámetros \\
    \bottomrule
  \end{tabularx}
  \caption{Módulos del editor visual}
\end{table}

\subsection{Flujo de creación de un componente}

\begin{enumerate}[label=\arabic*.]
  \item El usuario selecciona un tipo en la \textbf{Toolbar} (paleta izquierda).
  \item \texttt{circuitStore.addComponent()} crea un \texttt{CircuitComponent} con dos terminales y nodos eléctricos.
  \item React Flow renderiza un \texttt{ComponentNode} en la posición de la cuadrícula.
  \item El usuario conecta bornes con la herramienta \textbf{Cable}.
  \item \texttt{connectTerminals()} fusiona nodos eléctricos y registra un \texttt{Wire}.
\end{enumerate}

\subsection{Paneles de la interfaz}

\begin{table}[H]
  \centering
  \begin{tabularx}{\textwidth}{lX}
    \toprule
    \textbf{Panel} & \textbf{Archivo principal} \\
    \midrule
    Paleta de componentes & \texttt{components/toolbar/Toolbar.tsx} \\
    Propiedades y sondas & \texttt{components/properties/PropertiesPanel.tsx} \\
    Multímetro virtual & \texttt{components/multimeter/MultimeterDisplay.tsx} \\
    Osciloscopio & \texttt{components/graph/GraphPanel.tsx} \\
    Calculadora eléctrica & \texttt{components/calculator/CalculatorPage.tsx} \\
    Estado de simulación & \texttt{components/status/SimulationStatus.tsx} \\
    \bottomrule
  \end{tabularx}
  \caption{Paneles de la interfaz de usuario}
\end{table}

% ── 4. Motor de simulación ───────────────────────────────────────────
\section{Motor de simulación MNA}

\subsection{Fundamento matemático}

El análisis se basa en el método de \textbf{Análisis Nodal Modificado (MNA)}. Para cada instante de tiempo se construye y resuelve:

\begin{equation}
  \mathbf{G} \cdot \mathbf{x} = \mathbf{b}
\end{equation}

donde $\mathbf{G}$ es la matriz de admitancias del circuito, $\mathbf{b}$ el vector de fuentes y $\mathbf{x}$ el vector de incógnitas (voltajes de nodo y corrientes de ramas de fuentes de tensión).

Para el análisis \textbf{transitorio} se emplea integración \textbf{Euler hacia atrás} en capacitores e inductores. Los diodos y LEDs usan un modelo \textit{piecewise} con iteración de Newton (hasta 12 iteraciones por paso).

\subsection{Módulos del motor}

\begin{table}[H]
  \centering
  \small
  \begin{tabularx}{\textwidth}{lX}
    \toprule
    \textbf{Módulo} & \textbf{Responsabilidad} \\
    \midrule
    \texttt{SimulationEngine.ts} & Orquestador: sincroniza circuito, valida, avanza pasos \\
    \texttt{CircuitGraph.ts} & Grafo: terminales = vértices, elementos = aristas \\
    \texttt{ElementRegistry.ts} & Registro de modelos por tipo de componente \\
    \texttt{elements/index.ts} & \textit{Stamps} MNA (conductancias, fuentes) \\
    \texttt{MatrixBuilder.ts} & Construye $\mathbf{G}$ y $\mathbf{b}$ \\
    \texttt{TransientMNASolver.ts} & Resuelve MNA + transitorio \\
    \texttt{GaussianElimination.ts} & Eliminación gaussiana \\
    \texttt{CircuitValidator.ts} & Validación pre-simulación \\
    \bottomrule
  \end{tabularx}
  \caption{Módulos del motor TypeScript}
\end{table}

\subsection{Modelos de componentes}

\begin{table}[H]
  \centering
  \begin{tabular}{ll}
    \toprule
    \textbf{Componente} & \textbf{Modelo implementado} \\
    \midrule
    Resistencia & $G = 1/R$ \\
    Capacitor & Modelo companion (Euler) \\
    Inductor & Modelo companion (Euler) \\
    Batería & Fuente de tensión + ecuación de rama \\
    LED / Diodo & Piecewise: $R_{on}$, $R_{off}$, $V_f$ \\
    Interruptor & $R_{on}$ o $R_{off}$ según estado \\
    Amperímetro & $R = 0{,}1\,\Omega$ en serie \\
    Voltímetro & $R = 10\,\text{M}\Omega$ \\
    Transistor & Alta impedancia (placeholder) \\
    \bottomrule
  \end{tabular}
  \caption{Modelos eléctricos por componente}
\end{table}

% ── 5. Bucle de simulación ───────────────────────────────────────────
\section{Bucle de simulación en tiempo real}

Al pulsar \textbf{Iniciar simulación} ocurre la siguiente secuencia:

\begin{enumerate}[label=\arabic*.]
  \item \texttt{Toolbar} ejecuta \texttt{validateLocalCircuit()} --- si hay errores, muestra toast y no arranca.
  \item \texttt{toggleSimulation()} activa \texttt{simulationRunning = true}.
  \item \texttt{useSimulation} (hook) reinicia estado transitorio y datos del osciloscopio.
  \item Cada $\Delta t = 1/60$ s (60 Hz): \texttt{runLocalSimulationStep()} avanza un paso MNA.
  \item Se actualizan \texttt{simResults} (voltajes, corrientes) y \texttt{oscData} (series temporales).
  \item \texttt{ComponentNode}, \texttt{MultimeterDisplay} y \texttt{GraphPanel} se re-renderizan con los nuevos valores.
\end{enumerate}

\subsection{Instrumentación virtual}

\begin{itemize}
  \item \textbf{Multímetro:} al seleccionar un componente, muestra $V$, $I$ y potencia disipada. Casos especiales para amperímetro (corriente en serie), voltímetro (diferencia de potencial) y LED (caída directa y estado encendido/apagado).
  \item \textbf{Osciloscopio:} las sondas se añaden desde el panel de Propiedades con los botones \texttt{+V osc} y \texttt{+I osc}. Plotly.js renderiza las trazas con ventana deslizante de 10 s.
\end{itemize}

% ── 6. Circuito de demostración ──────────────────────────────────────
\section{Circuito de demostración}

El archivo \texttt{src/utils/loadDemo.ts} construye automáticamente un circuito en forma de \textbf{U} que incluye todos los tipos de componente del simulador:

\begin{verbatim}
Bat(+) → SW → Amm → R → D → LED → L → GND → Bat(−)
Paralelos: VM y Cap sobre LED; Pot sobre R; Trans sobre LED
\end{verbatim}

\textbf{Parámetros:} batería 9 V, resistencia 470 $\Omega$, LED $V_f = 2$ V, capacitor 47 $\mu$F, inductor 10 mH. Incluye cuatro sondas de osciloscopio preconfiguradas (corriente del amperímetro y voltajes de LED, capacitor e inductor).

Al iniciar la simulación, el LED se enciende en aproximadamente 1 segundo y la corriente de lazo es del orden de 12--14 mA.

% ── 7. Backend ───────────────────────────────────────────────────────
\section{Backend (FastAPI)}

El backend es \textbf{opcional} para el uso educativo diario, pero aporta:

\begin{itemize}
  \item \texttt{POST /api/simulate} --- simulación MNA en Python (referencia cruzada).
  \item \texttt{POST /api/simulate/validate} --- validación sin ejecutar.
  \item \texttt{POST /api/netlist} --- generación de netlist SPICE.
  \item \texttt{GET /api/health} --- comprobación de estado (\texttt{version: 0.1.2}).
\end{itemize}

\subsection{Estructura del backend}

\begin{verbatim}
backend/
├── main.py                    # Aplicación FastAPI + CORS
├── api/routes.py              # Endpoints REST
├── simulation/engine.py       # Motor MNA NumPy
├── validators/circuit_validator.py
├── models/circuit.py          # DTOs Pydantic
└── spice/builder.py           # Netlist SPICE
\end{verbatim}

% ── 8. Validaciones ──────────────────────────────────────────────────
\section{Validaciones del circuito}

Antes de simular, el validador comprueba:

\begin{itemize}
  \item Existencia de \textbf{un único nodo de tierra (GND)}.
  \item \textbf{Batería:} polo positivo conectado al circuito; polo negativo a GND.
  \item \textbf{Cortocircuitos} directos entre fuente y tierra.
  \item \textbf{Nodos flotantes} sin camino eléctrico a tierra.
  \item Cables duplicados entre los mismos terminales.
  \item Componentes sueltos (generan advertencia, no bloquean la simulación).
\end{itemize}

% ── 9. Despliegue ────────────────────────────────────────────────────
\section{Despliegue en la nube}

\begin{table}[H]
  \centering
  \begin{tabular}{lll}
    \toprule
    \textbf{Componente} & \textbf{Plataforma} & \textbf{Configuración} \\
    \midrule
    Frontend & Vercel & \texttt{vercel.json}, build Vite, salida \texttt{dist/} \\
    Backend & Render & \texttt{render.yaml}, \texttt{render-start.sh} \\
    Variable opcional & Vercel & \texttt{VITE\_API\_URL} = URL del backend \\
    \bottomrule
  \end{tabular}
  \caption{Despliegue del sistema}
\end{table}

La simulación del simulador funciona \textbf{sin backend}; Vercel sirve únicamente archivos estáticos del frontend.

% ── 10. Ejecución local ──────────────────────────────────────────────
\section{Ejecución local y pruebas}

\begin{lstlisting}
# Instalar dependencias
pnpm install

# Frontend (http://localhost:5174)
pnpm dev

# Backend opcional (http://localhost:8000)
python -m uvicorn backend.main:app --host 0.0.0.0 --port 8000

# Build de produccion
pnpm run build

# Test automatico del circuito demo
pnpm dlx tsx src/engine/demo/demoCircuitTest.ts
\end{lstlisting}

\subsection{Guion de demostración recomendado}

\begin{enumerate}[label=\arabic*.]
  \item Abrir la aplicación en el navegador.
  \item Pulsar \textbf{Demo} en la paleta izquierda.
  \item Pulsar \textbf{Iniciar simulación}.
  \item Observar el LED encendido, la corriente en el amperímetro y las cuatro trazas del osciloscopio.
  \item Seleccionar el LED y verificar el multímetro (caída directa, corriente, estado \textit{Encendido}).
\end{enumerate}

% ── 11. Equipo ───────────────────────────────────────────────────────
\section{Equipo de desarrollo}

\begin{table}[H]
  \centering
  \begin{tabular}{lll}
    \toprule
    \textbf{Integrante} & \textbf{Rama Git} & \textbf{Área principal} \\
    \midrule
    Luisa & \texttt{feature/luisa-backend} & Motor MNA, backend, validación \\
    Miguel & \texttt{feature/miguel-editor} & Editor React Flow, store, cables \\
    Josue & \texttt{feature/josue-ui} & UI/UX, paneles, documentación, despliegue \\
    \bottomrule
  \end{tabular}
  \caption{Distribución de responsabilidades}
\end{table}

% ── 12. Conclusiones ─────────────────────────────────────────────────
\section{Conclusiones}

Electro+ Lab demuestra la integración de:

\begin{itemize}
  \item \textbf{Ingeniería eléctrica:} análisis nodal modificado, modelos de componentes, instrumentación.
  \item \textbf{Desarrollo frontend:} React, TypeScript, estado reactivo, visualización con Plotly.
  \item \textbf{Algoritmos numéricos:} resolución de sistemas lineales, integración temporal, Newton para no linealidades.
  \item \textbf{Desarrollo backend:} API REST con FastAPI, validación y exportación SPICE.
  \item \textbf{DevOps:} despliegue en Vercel y Render con configuración reproducible.
\end{itemize}

El motor de simulación en el cliente garantiza respuesta inmediata sin latencia de red, lo que lo hace adecuado para uso en aula. El backend Python complementa con validación server-side y compatibilidad con herramientas SPICE estándar de la industria.

\vspace{1em}
\noindent\textbf{Documentación adicional en el repositorio:}
\texttt{docs/DOCUMENTACION\_PROFESOR.md},
\texttt{docs/arquitectura.md},
\texttt{docs/DEPLOY.md}.

\end{document}
